%%%%%
%
\section{R�gression lin�aire entre 2 param�tres}

\itemstyle{Op\'erateur~:} \operfmt{Reg\_lin}
\index{Reg\_lin}

%%%%%OBJET%%%%
\itemstyle{Objet~:}
Cet op�rateur calcule les coefficients de la r�gression lin�aire
par la m�thode des "moindres carr�s" et le param�tre estim�.

%%%%%ARGUMENTS%%%%
\bigskip
%%%%%ENTREES%%%%
\begin{tabarg}
\multicolumn{3}{l}{\titretab {Entr\'ees}}\\
\hline
$X$ & Vecteur, {\tt float} & Param�tre pr�dicteur\\
$Y$ & Vecteur, {\tt float} & Param�tre pr�dictande\\
\end{tabarg}

\medskip
%%%%%SORTIES%%%%
\begin{tabarg}
\multicolumn{3}{l}{\titretab {Sorties}}\\
\hline
$Z$ & Vecteur, {\tt float} & Param�tre estim�\\
\end{tabarg}


%%%%%FORMULE%%%%
\itemstyle{Formule}

%%Reg_lin

Les coefficients de la r�gression sont calcul�s par :
\begin{eqnarray*}
a &=& \frac{\overline {X.Y} - \overline{X}.\overline{Y} }
        {\overline {X^2} - \overline{X}.\overline{X}} \\
b &=& \overline{Y} - a. \overline{X} \\
\end{eqnarray*}

o� \[ \overline{X} = \frac{1}{N} \sum_{i=0}^{N-1} X[i] \]

Le param�tre estim� est calcul� par  :
\[
Z = a.X + b
\]



%%%%%COMMENTAIRE%%%%
\itemstyle{Commentaire~:}
Pour un gain de temps, cet op�rateur utilise
une formule de calcul des variances et covariance
qui ne n�cessite pas un calcul pr�alable de la moyenne.
Cette m�thode est sujette � des risques d'imprecision num�rique
lorsque le nombre d'�chantillons est grand (> 10$^{4}$),
et il est alors pr�f�rable d'utiliser l'op�rateur
\hyperref[OP-Reg_lin2]{\operfmt{Reg\_lin2}}.

%%%%%AUTEUR%%%%
\itemstyle{Auteur~:}
CNRM/GMEI/TRAMM

%%%%%REFERENCE%%%%
\itemstyle{R\'ef\'erences~:}
Vialar 1977 \cite{VIALAR-T3}.
